% conclusion.tex

\chapter{总结与展望}
\echapter{Conclusion and Future Work}

\section{本文工作总结}
\esection{Summary}

本文围绕时序数据预测任务，对轻量级多尺度模型 TimeMixer 进行了结构改进与实验分析。原始 TimeMixer 通过多尺度序列构建、趋势季节分解和跨尺度混合实现预测，具有较好的效率和结构优势。但从模型流程和任务特点看，原始模型仍存在平均池化下采样可能丢失高频信息、固定窗口分解难以适应复杂周期变化、通道独立策略忽略变量间相关性等问题。

针对上述问题，本文主要完成了三方面工作。第一，设计并实现了深度可分离卷积下采样模块，用可学习卷积替代固定平均池化，使模型在构建多尺度序列时具有一定的自适应特征选择能力。第二，设计并实现了多核门控分解模块，使用 $(3,11,37)$ 三个不同窗口并行提取趋势估计，再根据样本级全局均值特征生成门控权重并融合不同尺度结果。第三，在多尺度序列进入通道独立建模前加入低秩 MLP 形式的通道混合模块，使模型在保持轻量化的同时具备一定变量间信息交互能力。

实验部分，本文构造了 M0、M1、M2、M3、M12、M13、M23 和 M4 共 8 个模型变体，在 ETTh1、ETTh2、ETTm1、ETTm2、Weather、ECL、Traffic 和 Solar 8 个公开数据集上进行实验，预测长度包括 96、192、336 和 720。在固定随机种子 2021、实验重复次数为 1 的设置下，M2 在 64 个 MSE/MAE 评价项中取得 22 次最优，M0 取得 19 次，M4 取得 5 次；M1 和 M3 分别取得 3 次和 4 次最优。该结果说明，多核门控分解在本文实验中贡献最集中，而完整模型 M4 并未因包含全部模块而取得最多最优项。

此外，本文将内部变体结果与 PatchTST、Autoformer、DLinear、TimesNet、FEDformer 和 Crossformer 等模型的公开结果进行了平均结果比较。由于外部结果的训练环境和部分实验设置不由本文统一控制，该部分仅用于观察本文结果在常见模型中的相对位置。

总体来看，本文工作表明，改进 TimeMixer 的分解环节比单纯替换下采样或增加浅层通道混合更能集中改善本文实验中的预测结果。不同模块之间存在耦合关系，简单叠加并不一定带来更低误差；这一点也体现了完整消融实验对结构改进分析的必要性。

\section{不足与展望}
\esection{Limitations and Future Work}

本文仍存在一些不足。首先，深度可分离卷积下采样模块虽然增强了下采样的学习能力，但 M1 在 64 个评价项中只取得 3 次最优。对于较平滑的数据，平均池化的平滑作用可能反而更有利；对于噪声较大的数据，可学习卷积也可能保留不必要的扰动。后续可以考虑在卷积下采样中加入更明确的平滑约束，或者根据数据特征自适应选择下采样方式。

其次，通道混合模块目前只在多尺度序列进入通道独立建模前进行一次浅层变量交互，表达能力有限。对于变量数较多且变量关系复杂的数据集，仅使用一个低秩 MLP 可能不足以充分建模变量之间的依赖关系。后续可以尝试分层通道混合、注意力式变量选择或图结构建模等方法，但也需要注意控制模型复杂度。

再次，多核门控分解模块虽然在本文设置下取得最多最优项，但窗口大小仍采用固定设置 $(3,11,37)$，尚未系统分析不同窗口组合对性能的影响。后续可以围绕短窗口、中窗口和长窗口组合开展参数敏感性实验，例如比较 $(3,7,15)$、$(3,11,37)$ 和 $(5,25,49)$ 等设置，以进一步说明窗口大小与数据周期结构之间的关系。

此外，本文各实验设置仅在固定随机种子下运行一次，尚未统计多随机种子下的均值和方差，因此对模型稳定性的判断主要基于当前实验设置，后续可补充多次重复实验以提高结论可靠性。

最后，本文实验主要围绕公开数据集和消融分析展开，对模型参数量、训练时间、推理耗时和显存占用等效率指标的系统比较仍不充分。由于 TimeMixer 本身强调轻量化，后续若要更全面评价改进模型，应在统一环境下补充效率指标对比。同时，本文外部模型对比主要基于原论文公开结果，后续如果条件允许，可以在统一实验条件下重新复现更多 baseline。除此之外，还可以在更多真实业务数据上验证模型泛化能力，使研究结论更接近实际应用需求。
